% Part: lambda-calculus
% Chapter: introduction
% Section: term-revisited

\documentclass[../../../include/open-logic-section]{subfiles}

\begin{document}

\olfileid{lam}{syn}{tr}
\olsection{ٹرم بطور $\alpha$-ہم ارزیتی طبقے}

اگّوں اسیں ٹرماں نوں $\alpha$-ہم ارزی دے حساب نال ویکھاں گے۔
ایہدا مطلب اے کہ جدوں اسیں کوئی ٹرم لکھدے آں، تاں ساڈی مراد
اوہدا $\alpha$-ہم ارزیتی طبقہ اے۔ مثلاً، $\lambd[a][\lambd[b][a c]]$
اوہناں سارے ٹرماں دا سیٹ دسدا اے جیہڑے اوہدے نال $\alpha$-ہم ارز نیں،
جیویں $\lambd[a][\lambd[b][a
    c]]$، $\lambd[b][\lambd[a][b c]]$ وغیرہ۔

پچھلے حصیاں وچ $N, Q$ ورگے حرف اک ٹرم دسدے سن۔ اگّوں
اسیں اوہی حرف طبقہ دسن لئی ورتاں گے، تے ہِٹھاں ساڈا مطالعہ
ٹرماں دی تھاں ایہناں طبقیاں بارے ہووے گا۔ $x, y$ ورگے حرف
ہلے وی متغیر دسّن گے۔

اسیں $\rep{M}$ دی لکھت طبقہ~$M$ دے کسے وی !!{element}
لئی ورتاں گے، تے جے اک توں ودھ چاہیدے ہون تاں $\rep{M}[0], \rep{M}[1]$ وغیرہ۔

گل نوں سوکھا رکھن لئی ٹرماں والی لکھت فیر ورتدے آں۔
طبقیاں لئی تعریف ہِٹھ دِتّی اے۔
\begin{defn}
  \begin{enumerate}
  \item $\lambd[x][N]$ اوہ طبقہ اے جیہدے وچ
    $\lambd[x][\rep{N}]$ شامل اے۔
  \item $PQ$ اوہ طبقہ اے جیہدے وچ $\rep{P}\rep{Q}$ شامل اے۔
  \end{enumerate}
\end{defn}

ویکھنا اوکھا نہیں کہ ایہ ٹھیک طرح معین نیں، کیوں جے
$\alpha$-روپ بدلن دا تعلق موافق اے۔

\begin{defn} \ollabel{def:fv}
  $\alpha$-ہم ارزیتی طبقہ~$M$ دے \emph{آزاد متغیر}، یعنی
  $FV(M)$، دی تعریف $FV(\rep{M})$ اے۔
\end{defn}

ایہ ٹھیک طرح معین اے، کیوں جے $FV(\rep{M}[0]) =
FV(\rep{M}[1])$، جیویں \olref[alp]{thm:fv} وچ دکھایا اے۔

طبقیاں وچ بدل لئی وی اوہی لکھت ورتدے آں۔
\begin{defn} \ollabel{def:sub}
  $R$ نوں $y$ دی تھاں $M$ وچ رکھن والا \emph{بدل}، یعنی
  $\Subst{M}{R}{y}$، دی تعریف $\Subst{\rep{M}}{\rep{R}}{y}$ اے،
  کسے وی ایہو جیہے $\rep{M}$ تے $\rep{R}$ لئی جیہڑے بدل نوں
  معین کردے ہون۔
\end{defn}

ایہ وی ٹھیک طرح معین اے، جیویں \olref[alp]{cor:sub} وچ دکھایا اے۔

دھیان دیو کہ ایس تعریف نال ساڈا استدلال کِنّا سوکھا ہو جاندا
اے۔ ایہدی اک مثال ہِٹھ اے۔
\begin{align}
  \Subst{\lambd[x][x]}{y}{x} & =\ollabel{eq:1}\\
  &= \Subst{\lambd[z][z]}{y}{x} \ollabel{eq:2}\\
  &= \lambd[z][\Subst{z}{y}{x}] \\
  &= \lambd[z][z]
\end{align}

جے \olref{eq:1} نوں ہلے وی ٹرماں والا بدل منّیئے، تاں ایہ
غیر معین اے۔ پر جیویں پہلاں آکھیا اے، ہُن اسیں اینوں طبقیاں
والا بدل سمجھدے آں۔ ایس لئی \olref{eq:2} ممکن اے۔
$\lambd[x][x]$ دی تھاں $\lambd[z][z]$ رکھ سکدے آں، کیوں جے
دوویں اکّو طبقہ وچ نیں۔

ایسے کارن اگّوں اسیں منّاں گے کہ ساڈے چُنے ہوئے نمائندے بدل
لئی لوڑ دیاں شرطاں ہمیشہ پوریاں کردے نیں۔ مثلاً، جدوں
$\Subst{\lambd[x][N]}{R}{y}$ ویکھیے، تاں منّاں گے کہ نمائندہ
$\lambd[x][N]$ ایویں چُنیا اے کہ $x \neq y$ تے $x \notin FV(R)$۔

کیوں جے $\lambd[x][x]$ نوں «طبقہ» آکھنا ذرا اوپرا لگدا اے،
اگّوں ایہناں نوں $\Lambda$-ٹرم آکھدے آں، یا ایس حصے دے باقی
متن وچ بس «ٹرم»۔ ایس طرح ایہ پہلاں توں جانے ہوئے $\lambda$-ٹرماں
توں وکھرے رہن گے۔

\begin{editorial}
  ہلے اسیں ٹرماں نوں نہیں چھڈ سکدے۔ $\Lambda$-ٹرماں دی ساری
  تعریف $\lambda$-ٹرماں اُتے ٹکی اے، تے اسیں $\Lambda$-ٹرماں
  اُتے فنکشن متعین کرن دا کوئی طریقہ ہلے نہیں دِتّا۔ ایس دا
  مطلب اے کہ ایہو جیہا ہر فنکشن پہلاں $\lambda$-ٹرماں اُتے
  متعین کرنا پینا اے، تے فیر $\Lambda$-ٹرماں اُتے «منتقل» کرنا
  پینا اے، جیویں اسیں بدلاں لئی کیتا سی۔ فیر وی اسیں منّ کے
  چلدے آں کہ پڑھن والا بدیہی طور اُتے سمجھ سکدا اے کہ
  $\Lambda$-ٹرماں اُتے فنکشن کیویں متعین ہو سکدے نیں۔
\end{editorial}
\end{document}
